\chapter{فهم المرصد ورحلة العمل}
\label{ch:01-overview}
\section{لماذا نستخدم المرصد؟}
\label{sec:observatory-purpose}

حتى نعرف واقع الصناعة، نحتاج إلى تحديد المنشآت التي ندرسها، وجمع بياناتها عن فترة محددة، ومراجعتها قبل قراءة النتائج. وإذا جُمعت الإجابات في ملفات متفرقة، يصعب معرفة من أجاب، وما الذي لم يُستكمل، وهل تتعلق الأرقام بالفترة نفسها.

يجمع كيان هذه الأعمال في رحلة مترابطة: يعرّف السجل الصناعي المنشأة، وتحدد خطة التكليف العمل المطلوب ومن ينفذه، وتجمع النماذج الإجابات، ثم تساعد صفحات المتابعة والمؤشرات على قراءة ما سُجّل. لا تعني كثرة الإجابات أنها مكتملة أو معتمدة؛ فجمع البيانات ومراجعتها مرحلتان مختلفتان.

\begin{rcnote}[تطبيق المثال: مصنع الأفق للصناعات الغذائية]
نريد معرفة عدد العاملين وأيام توقف التشغيل في مصنع الأفق بريف دمشق خلال أيلول 2026. سنستخدم هذا المصنع الافتراضي طوال الدليل، ونبدأ بسؤال واضح: ماذا حدث خلال هذه الفترة، وما الاحتياج الذي يستدعي المتابعة؟
\end{rcnote}

في مثالنا، تتولى سارة إعداد المسح وخطته، ويجمع عمر البيانات، وتراجع ليلى الإجابات. هذه مسؤوليات افتراضية لتوضيح العمل، وليست أسماء صلاحيات جاهزة في النظام. سنحدد طريقة منح الصلاحيات عند شرح الأدوار.

\section{ما الفرق بين السجل الصناعي والمرصد؟}
\label{sec:observatory-register}

السجل الصناعي يجمع ملف المنشأة: هويتها، والجهة القانونية المرتبطة بها، وقيدها الصناعي وأصولها. أما المرصد فيجمع بيانات المسح ومتابعته ومؤشراته، ومنها بيانات تتعلق بفترة أو جولة محددة. وقد تتغير بيانات ملف المنشأة أيضاً؛ فالفرق هو الغرض من البيانات، وليس كونها ثابتة إلى الأبد.

\begin{table}[H]
\centering
\caption{سؤالان مختلفان عن المنشأة نفسها}
\label{tab:observatory-register}
\begin{tabularx}{\linewidth}{Y Y}
\rcHeaderRow
\thh{السجل الصناعي} & \thh{المرصد الصناعي}\\
ما المنشأة التي نتحدث عنها؟ & ما بياناتها في فترة المسح؟\\
من مالكها وما موقعها وما أصولها؟ & ما الذي جُمع، وما الذي يحتاج إلى استكمال أو مراجعة؟\\
يساعد على تحديد ملف المنشأة المقصود. & يساعد على متابعة المسح وقراءة نتائجه.\\
\end{tabularx}
\end{table}

\begin{rcnote}[تطبيق المثال]
اسم «مصنع الأفق للصناعات الغذائية» وموقعه في ريف دمشق يعرّفان المنشأة في مثالنا. أما «24 عاملاً و3 أيام توقف خلال أيلول 2026» فهي بيانات المسح التي سنجمعها عنها؛ لا يكفي حفظ العدد دون معرفة المنشأة والفترة اللتين يخصهما.
\end{rcnote}

\textbf{تحقق من فهمك:} إذا بدأنا مسح تشرين الأول للمصنع نفسه، فهل أصبح مصنعاً جديداً؟ لا؛ نتعامل مع المنشأة نفسها، لكننا نجمع بيانات فترة أخرى. وسنشرح إعداد الجولة وربطها بالسجلات في موضعه، دون افتراض أن تغيير التاريخ وحده ينشئ جولة.

\section{كيف تتكامل النماذج والحزمة والخطة؟}
\label{sec:observatory-building-blocks}

\textbf{النموذج} يحدد الأسئلة التي نريد إجاباتها والقواعد التي تضبط إدخالها. و\textbf{حزمة النماذج} تجمع نماذج متعددة لمتابعتها ضمن عمل واحد مرتبط بسجل؛ لكل وحدة فيها شروط فتح واكتمال بحسب إعداد الحزمة.

\textbf{خطة التكليف} تحدد السجلات المستهدفة، والمرشحين للعمل، وطريقة توزيعه ومواعيده. أما \textbf{التكليف} فهو العمل الناتج والمسنَد لتنفيذ الحزمة على سجل معين. و\textbf{الجولة} تميز دورة تنفيذ الخطة، مع الفترة المحددة لها بحسب الإعداد.

\begin{rcnote}[تطبيق المثال]
سنُعدّ نموذج «بيانات التشغيل» لعدد العاملين وأيام التوقف، ونموذجاً مساعداً صغيراً باسم «متابعة الاحتياج». سنجمعهما في حزمة «مسح تشغيل الأفق»، ثم نستخدم الحزمة في خطة «متابعة تشغيل الأفق» لاستهداف المصنع وإسناد العمل إلى عمر بعد تجهيز قواعد الإسناد.
\end{rcnote}

لا يكفي إنشاء النموذج ليصل العمل إلى عمر؛ يجب إعداد الحزمة والخطة والمرشحين ثم التحقق من الإسناد. كما أن جمع نموذجين في حزمة لا يعني انتقال الإجابات بينهما تلقائياً؛ سنشرح الربط الفعلي وشروطه عندما نبني المثال.

في هذا الدليل نبني نموذجاً واحداً بالتفصيل، ونضيف النموذج المساعد لتوضيح الربط. لا تحتاج الآن إلى تعلّم النماذج الثلاثة عشر الموجودة في المسح الحالي.

\section{اتبع رحلة المثال من التجهيز إلى النتيجة}
\label{sec:observatory-journey}

نبدأ بتجهيز ما يحتاج إليه المسح، ثم نوزع العمل ونجمع الإجابات ونراجعها. لا تُغني مرحلة عن الأخرى: وجود الخطة لا يثبت إسناد كل منشأة، وإرسال الإجابة لا يثبت اعتمادها.

\begin{rcnote}[تطبيق المثال]
سننفذ الرحلة التالية تدريجياً في الفصول القادمة، بالأسماء والقيم نفسها. هذا مخطط للمثال الذي سنطبقه، وليس تقريراً عن سجلات أُنشئت أو إجراءات اختُبرت بالفعل.
\end{rcnote}

\begin{figure}[H]
\centering
\begin{tikzpicture}[node distance=0.4cm,stepBox/.append style={minimum width=6cm,text width=5.6cm,minimum height=0.7cm}]
  \node[stepBox] (s0) {\RL{نجهّز ملف مصنع الأفق}\\\RL{وفريق العمل}};
  \node[stepBox,below=of s0] (s1) {\RL{نبني نموذج التشغيل}\\\RL{ونجهّز الحزمة}};
  \node[stepBox,below=of s1] (s2) {\RL{تجهّز سارة خطة التكليف}\\\RL{وتتحقق من إسناد العمل}};
  \node[stepBox,below=of s2] (s3) {\RL{يعبّئ عمر النماذج}\\\RL{ويرسل الإجابات}};
  \node[stepBox,below=of s3] (s4) {\RL{تراجع ليلى البيانات}\\\RL{ثم نقرأ النتائج}};
  \draw[flowsTo] (s0) -- (s1);
  \draw[flowsTo] (s1) -- (s2);
  \draw[flowsTo] (s2) -- (s3);
  \draw[flowsTo] (s3) -- (s4);
\end{tikzpicture}
\caption{رحلة العمل المقترحة في مثال مصنع الأفق}
\label{fig:observatory-journey}
\shotnote{يربط المخطط تجهيز ملف المصنع والنماذج بالعمل الذي سينفذه عمر وتراجعه ليلى. يوضح ترتيب المثال ومسؤولية كل مرحلة؛ ولا يعني أن الإسناد أو الإرسال أو الاعتماد يحدث تلقائياً بمجرد الانتقال إلى المرحلة التالية.}
\end{figure}

\begin{enumerate}
\item نحدد مصنع الأفق والفريق الذي سيتولى مسحه. نتحقق من أن العمل يتعلق بالمنشأة المقصودة، وأن المسؤولين معروفون.
\item نبني «بيانات التشغيل» ثم نضيف «متابعة الاحتياج» إلى الحزمة. نختبر الأسئلة والقواعد قبل استخدام الحزمة في الخطة.
\item تجهّز سارة الخطة لجولة أيلول 2026 وتتحقق من السجلات والمرشحين. بعد الإطلاق يجب التأكد من وجود تكليف المصنع وإسناده إلى عمر؛ وإذا لم يُسند، تُراجع شروط الترشيح والسعة بدلاً من افتراض نجاح التوزيع.
\item يسجل عمر 24 عاملاً و3 أيام توقف، ويذكر انقطاع الكهرباء سبباً للتوقف، ثم يرسل البيانات وفق المسار الذي سنعدّه. النتيجة المطلوبة هنا إجابات قابلة للمراجعة، لا اعتماد نهائي.
\item تراجع ليلى الأرقام والفترة وسبب التوقف. نعالج الملاحظات وفق مسار المراجعة المهيأ، ثم نقرأ النتائج مع الانتباه إلى حالة البيانات التي تعرضها اللوحة.
\end{enumerate}

\section{كيف نعرف ما أُنجز وما بقي؟}
\label{sec:observatory-follow-up}

تختلف متابعة التنفيذ عن قراءة النتائج. في التنفيذ نسأل: هل أُسند العمل؟ هل بدأ العضو؟ وما النماذج التي تحتاج إلى استكمال أو مراجعة؟ أما عند قراءة النتائج فنسأل: ماذا تقول البيانات، وعن أي فترة، وما حالتها؟

قد يكون التكليف قد بدأ بينما بعض وحداته لم تكتمل. وقد تكون الإجابة مُرسلة وتنتظر المراجعة. لذلك لا نستخدم «بدأ» و«أرسل» و«اعتمد» بالمعنى نفسه، ولا نفترض أن اكتمال نموذج واحد يعني اكتمال الحزمة كلها؛ يعتمد ذلك على إعدادها.

\begin{rcnote}[تطبيق المثال]
إذا أرسل عمر «بيانات التشغيل» وبقي «متابعة الاحتياج» مطلوباً، فما زال هناك عمل لم يُستكمل في مثالنا. وإذا طلبت ليلى توضيح سبب التوقف، فالمطلوب متابعة هذه الملاحظة، لا الاكتفاء بوجود الرقم «3» في الإجابة.
\end{rcnote}

تعرض الخطة التي راجعناها في المرصد أعداداً منفصلة للعمل المسند والمبدوء والمكتمل والمعتمد والمتأخر. هذه الفروق تساعد المسؤول على تحديد ما يحتاج إلى متابعة؛ أما معاني الحالات التفصيلية وشروط انتقالها فسنشرحها مع الإعداد الفعلي.

\section{طبّق الفكرة قبل أن نبدأ الإعداد}
\label{sec:observatory-exercise}

قبل استخدام الأزرار، حدّد الشيء الذي ستنشئه ولماذا تحتاج إليه. التمرين التالي يثبت صورة العمل؛ لا يتطلب إنشاء سجلات في كيان الآن. سنحوّله إلى تطبيق داخل النظام في الفصول التنفيذية.

\begin{rcnote}[تطبيق المثال]
اكتب بطاقة قصيرة لمصنع الأفق تتضمن المنشأة، وفترة المسح، والأسئلة، والحزمة، ومن يجهز العمل وينفذه ويراجعه. قارن بطاقتك بالقيم التالية؛ يجب أن تعرف أي عنصر يحدد المنشأة، وأي عنصر يحدد البيانات المطلوبة، ومن سينفذ العمل.
\end{rcnote}

\begin{table}[H]
\centering
\caption{بطاقة العمل التي سنبني عليها}
\label{tab:observatory-example-card}
\begin{tabularx}{\linewidth}{R{3.2cm} Y}
\rcHeaderRow
\thh{العنصر} & \thh{القيمة في المثال}\\
المنشأة & مصنع الأفق للصناعات الغذائية، ريف دمشق.\\
فترة المسح & أيلول 2026.\\
النموذج الأساسي & بيانات التشغيل: 24 عاملاً، و3 أيام توقف، والسبب انقطاع الكهرباء.\\
النموذج المساعد & متابعة الاحتياج: تحديد الاحتياج الذي يستدعي المتابعة، دون افتراض أنه طلب خدمة أو موافقة على دعم.\\
الحزمة & مسح تشغيل الأفق.\\
الخطة & متابعة تشغيل الأفق.\\
المسؤوليات & سارة تجهّز، وعمر يجمع البيانات، وليلى تراجع.\\
\end{tabularx}
\end{table}

\begin{rcchecklist}[تحقق من فهمك]
\begin{checklist}
\item أستطيع التمييز بين ملف مصنع الأفق وبيانات المسح الخاصة بأيلول.
\item أعرف أن النموذج يحدد الأسئلة، وأن الخطة تحدد نطاق التكليف وتوزيعه.
\item أستطيع شرح الفرق بين جمع النماذج في حزمة وانتقال البيانات بينها.
\item أعرف من يجهّز المثال ومن ينفذه ومن يراجعه، وأن صلاحياتهم تحتاج إلى إعداد.
\item لا أعدّ إرسال عمر للإجابة اعتماداً لها أو اكتمالاً لكل العمل.
\end{checklist}
\end{rcchecklist}

بعد مراجعة هذا الفصل واعتماده، ننتقل إلى الفصل~\ref{ch:02-roles} لتحديد الأعضاء والأدوار ونقطة البداية لكل مشارك. ستبقى المنشأة والأسماء والحزمة والخطة نفسها معنا في بقية الرحلة.

